\clearpage
\chapter*{第二部分结语}
\addcontentsline{toc}{chapter}{第二部分结语}
\label{part:02-closing}

\begin{center}
  \includegraphics[width=\linewidth,height=.43\textheight,keepaspectratio]{assets/part-art/part-2-closing.png}
  \par\vspace{4pt}
  {\scriptsize\color{Muted}取意于康斯太勃尔《水闸》\par}
  \BookEpigraph{part2}
\end{center}

Alice 的购买已经跨过几次进程退出。平台扣去 3000 分，购买服务一度还在等待付款；后来订单与通知一起保存，消息也可能发过两次。同一 EV1 先后到达履约处理器三次，后两次都找到原来的 F1。把这段历史按系统分开看，每个系统都保存着自己已经确认的事实，恢复程序沿这些事实接续下一段工作。

第一处变化发生在 A。购买服务先保存 P1、C1 与固定平台身份，后来工作者才能用原输入继续付款。E 已在平台提交、R 尚未在购买库完成时，本地等待准确表达了证据尚缺的阶段。查询、回调与对账逐步补足这份证据，R 再把原结果和唯一订单共同保存。超时留下的两个分支由原身份关联，服务可以继续查询和恢复，而无须从超时本身猜出平台是否已经扣款。

T1 把持续运行的责任留在了数据库。进程失联以后，扫描器根据状态与到期条件重新发现它；新的领取取得新的代次，调度入口据此拒绝旧工作者的回写。与此同时，一份可信的旧平台响应依然可能说明 X1 已经存在。授权谁安排下一轮工作，与确认付款发生了什么，分别在自己的入口核验。到了独立资源，接收者还要把代次检查、水位推进和实际修改共同仲裁，才能拒绝迟到旧写。同代次不同业务进度的顺序，则由另外保存的业务版本解释。

通知把相似的问题带到新的边界。R 可以原子保存付款结论与 EV1，代理稍后独立受理消息。发布确认丢失以后，PA2 可能把同一 EV1 再送一次，因此消费者继续携带原事件身份。Q 用 I 恢复一次接收，用 B 约束初次履约，把真实 F 与本地结果、回报责任一起提交。另一份合法购买取得另一 O，两个不同事件也可能报告同一个 O；两层身份让这两种情况得到各自正确的处理。

\clearpage
\CTenFigure{part02-responsibility}{上游先保存待办与下游身份，后续接收者提交自己的效果及继续依据}{fig:part02-responsibility}

这些步骤共有一项构造方法。先找到当前系统能够原子保存的那组事实，把下一段工作的原输入、稳定身份和待办放进这次提交。接收者取得工作以后，在自己能控制效果的位置仲裁重复、核对输入并保存结果。发送者缺少确认时，从原记录继续；收到足够证据时，结束本段待办。每次转交都留下可查的位置，进程和连接的寿命便可以短于整项业务的寿命。

这条方法也限定了能够作出的承诺。外部平台缺少固定身份保护时，本地唯一行数无法限制平台扣款数。代理确认以后若允许无证据删除消息，消费者可能失去继续输入。仓储请求已经发出却无法查明结果时，持久核查保住原映射和处置责任，发货结局则等待相应证据。安全保证逐段依赖具体仲裁，最终完成逐段依赖记录可读、实际调度、网络机会与有依据的处置。

本地回滚的作用域始终清楚。R 回滚撤销购买库内尚未提交的订单、结果和事件，平台 E 继续保留已完成的收款。Q 回滚撤销履约库内的候选，代理中的副本继续按自己的规则存放。后来需要退款、补发或更正时，它们各自取得业务授权、身份和恢复路径。已经确认的付款与履约成为这些后续动作的依据。

到这里，每一段恢复都用到了已提交记录的延续性：P/C 的固定映射、任务代次、资源高水位、EV 的原快照，以及 I/B/F 的绑定。这个条件在前面采用的故障模型中成立，接下来的系统变化会直接检验它。若新主库缺少先前确认的购买，旧请求会再次找到空位置；若去重结果已经回收，迟到副本可能失去拒绝依据；若备份把任务代次恢复到昨天，旧消息和新领取又会共享过去用过的数字。

第三部分从这些具体历史继续：复制怎样保住已承诺的事实，记录在什么条件下才可遗忘，以及恢复旧历史以后，服务要取得哪些外部证据才能重新开放入口。一次购买的身份继续沿用，变化的是恢复者还能依赖哪些保存下来的历史。
